\chapter{Introduction}
\label{cha:introduction}

From its beginnings in the 1990s, the Web has become an important
application platform in less than one and a half decades. Industry
surveys indicate that an increasing number of companies develop
distributed Web applications as a core part of their
business~\cite{kondo2009cost,EYPaaSSurvey2012}. However, many of those
applications are not written with distributed systems middleware to
handle communication between services. Instead, developers improvise
their own communication interfaces. This is error-prone and
complicates integration~\cite{Nadiminti}.

Distributed applications are especially prone to software
failures. Recent studies~\cite{Ballou2008,Tassey2002} show a relative
increase in code problems, hidden dependencies, and accidental
complexity when integrating multiple systems. Different systems may
provide different network interfaces and developers need to understand
their utilization and implicit assumptions. Different systems also run
different frameworks and programming languages, so that developers
need help to associate observed network communication with code
locations. This additional complexity makes it harder for developers
to reason about the system as a whole, thus, they are more likely to
introduce defects~\cite{zeller2009programs}. Debugging the resulting
failures is hard, and we have identified four primary problems that
developers face with current approaches and tools.

First, reliably reproducing failures is hard. Since distributed
systems execute in parallel, a global event order is difficult to
establish and may vary non-deterministically across
executions~\cite{lamport1978time}. Such variance depends on the
physical connections between the servers. The development setup will
differ from deployment and, thus, network delays and probabilities for
non-deterministic behavior also differ. Consequently, developers may
not be able to trigger some failures during development that occur in
deployment, or such failures may occur so rarely that it is impossible
to debug them effectively, because developers cannot reliably
reproduce them.

Second, developers cannot easily distinguish between causes and
effects. Many developers rely on systematic logging to understand
failures that only occur during
deployment~\cite{zeller2009programs}. However, isolated event logs are
poorly suited to create a global view on the events in a distributed
application. The ability to reason about global event order is
necessary for humans to understand causality in a
system~\cite{lamport1978time}. Understanding the cause of a failure is
essential for fixing it correctly.

% Consider a search engine. For a small scale engine, developers can
% control data storage, transactions, and database connections. On the
% Internet, the primary communication protocol is \ac{http} which limits
% in the amount of control developers have over their systems'
% connections. Non-atomic data updates usually cannot be wrapped in
% transactions. Storage services on the Web cannot provide transactions
% to thousands of users in parallel. Parallel \ac{http} request are each
% stateless and provide no guaranteed ordering. Thus, developers can
% rely less on the underlying infrastructure and require debugging tools
% work in a distributed setting.

Third, developers cannot easily inspect execution on servers in the
system. With symbolic debuggers, programmers inspect execution state
and behavior to find anomalies that may have caused a
failure. Symbolic debuggers enable them to stop execution and
interactively inspect the program's state. However, in a distributed
system, developers need to inspect state across multiple servers,
which execute in parallel. Symbolic debuggers are tailored towards
systems under developer control, but not all parts of a distributed
system provide that kind of control. Stopping servers with a debugger
to inspect the execution state is not feasible during deployment and
external services may not provide any debugging mechanism.

Finally, different frameworks and programming languages on several
servers force developers to apply multiple debugging
tools. Consequently, developers cannot debug and inspect the
distributed system as one. Instead, they have to treat each system
separately, because current debugging tools cannot connect and
synchronize with others. Support for debuggers that work across
multiple servers and languages is virtually non existent.

Due to these problems, debugging failures in distributed systems,
especially non-deterministic failures, is a time-consuming task. It
usually consists of aggregating various forms of log files, debug
output, and \ac{api} documentation for the employed services, and then
debugging on an equivalent, local system to infer possible
problems. Symbolic debuggers cannot be used in distributed
applications, because they only support independent analysis of
individual processes and require full control over those.

Tool support for network observation, log aggregation and debugging
across different systems is virtually non-existent. Yet, to improve
programmer productivity, providing adequate debugging tools is a key
ingredient, as without them, debugging is a manual, time-consuming
task~\cite{Dinning1990}.

We argue that there is a need for debugging tools that target
distributed, heterogeneous systems which are not entirely under
developer control. Architectural patterns for the these kinds of
systems, such as \ac{soa} and \ac{rest}; and hosting services, such as
\ac{paas} and \ac{iaas}, are relatively new phenomena---the first
comprehensive description of \ac{soa} was in 2005~\cite{Erl2005}. For
that reason, much of the work in this area is still not aimed at the
specific challenges of debugging distributed applications. Developers
of these systems need tools that provide a means to record and replay
events consistently across multiple servers. The tools should
furthermore guide developers towards anomalous differences across
multiple executions, so developers can reason about causality in the
system. Finally, such a tool should allow developers to cope with
multiple frameworks and languages, as they are common in distributed
applications.

\section{Approach}
\label{sec:approach}

In this work we present an approach for integrated debugging in
distributed applications, targeted specifically at Web development,
but extensible to other kinds of applications with similar
problems. Our approach---Replay-driven Fault Navigation---combines
state of the art techniques in the area of record and replay to
eliminate of non-determinism, distributed program comprehension, and
statistical debugging.

To allow debugging of non-deterministic failures, we present a tracing
approach for unified monitoring of systems. We continuously monitor
the live system and record enough data to allow reproducing relevant
program state at a later time. Our monitoring imposes little overhead
so it does not disturb execution unduly and can stay enabled during
deployment. We analyze the recorded data and correlate it to failures
to determine causes and effects. Network schedules are treated as
sequences of send and receive events between pairs of servers. Our
heuristics assign cause probabilities to sub-sequences of the recorded
network schedules, to guide developers during debugging. Developers
can replay such event sequences with our network scheduler and
iteratively refine recorded execution data for inspection. The
scheduler is an intermediate layer that controls how network events
pass between the application code and the framework. It blocks or
stalls events to change the order in which they pass in and out of the
application to match the recorded sequence. We correlate the network
data with source code and present both in a unified debugging tool, so
developers can debug at either level of abstraction and correlate
source code across servers and languages.

The advantages of our approach over traditional stop-the-world
symbolic debuggers are: the impact on the runtime behavior is kept at
a minimum, the technique can be applied to the live system,
eliminating the need for a test setup, and the debugging tools can
help developers to logically connect source code locations across
different servers.

To our knowledge, the presented approach is unique in that it
explicitly tries to support debugging across different languages and
system with one integrated approach. Work on understanding and
debugging distributed systems has been on-going for some
time~\cite{Parker1983} and the state of the art has solutions for each
part of Replay-driven Fault Navigation, but no integrated solution
exists that we are aware of.

\section{Outline}
\label{sec:outline}

This work is structured as follows: in \autoref{sec:requ-based-distr},
we discuss the problem domain on a small example. We highlight
challenges faced by developers when working with current debugging
techniques. In \autoref{cha:solution-design}, we present our four-step
approach to those challenges, and how that approach could work in
different languages and deployment scenarios. In
\autoref{cha:implementation}, we discuss an implementation of
Replay-driven Fault Navigation in Smalltalk. Finally in
\autoref{cha:evaluation}, we use the initial example and a larger
application adapted from an industrial system to evaluate our tool in
comparison with contemporary techniques. In
\autoref{cha:related-work}, we discuss related work and present our
conclusions in \autoref{cha:conclusion}.

%%% Local Variables: 
%%% mode: latex
%%% TeX-master: "thesis"
%%% End: 
